\documentclass[11pt,a4paper]{article}

% ── Packages ─────────────────────────────────────────────────────
\usepackage[utf8]{inputenc}
\usepackage[T1]{fontenc}
\usepackage{amsmath,amssymb}
\usepackage{booktabs}
\usepackage{graphicx}
\usepackage{xcolor}
\usepackage{hyperref}
\usepackage{enumitem}
\usepackage[margin=2.5cm]{geometry}

\hypersetup{
    colorlinks=true,
    linkcolor=blue!60!black,
    citecolor=blue!40!black,
    urlcolor=blue!60!black,
}

% ── Title ────────────────────────────────────────────────────────
\title{Response to Reviewers}
\author{
    \textbf{Manuscript:} CrossEcho: Predictive Modelling of Crosstalk Side-Channel\\
    Leakage on Multi-Tenant Superconducting QPUs\\[4pt]
    \textbf{Journal:} Quantum Engineering (Wiley)\\[4pt]
    \textbf{Date:} 8 August 2026
}
\date{}

\begin{document}
\maketitle

\noindent
We thank both reviewers for their careful reading and constructive feedback. We have
addressed all comments in detail below. Each reviewer comment is reproduced in
\textbf{bold italics}, followed by our response and a description of the corresponding
changes to the manuscript. All section, equation, figure, and table numbers refer to
the revised manuscript.

% ===========================================================================
\section*{Reviewer 1}
% ===========================================================================

% ── R1C1 ─────────────────────────────────────────────────────────
\subsection*{Comment 1}

\noindent\textbf{\textit{``Given that the readout-only control is a key methodology in your protocol,
does it mean measurement only (or measurement-induced crosstalk)? How to realize it
in practice, by rule or some technics?''}}

\medskip\noindent\textbf{Response:}

We thank the reviewer for this important clarification question. The readout-only control
is a \textbf{control experiment}, not part of the crosstalk attack or detection protocol
itself. Its purpose is to isolate and quantify measurement-induced (readout) crosstalk so
that it can be disambiguated from gate-induced ZZ crosstalk. The key idea is
straightforward: if the victim circuit has zero gates, any non-zero TVD measured in the
attacker's statistics cannot arise from gate-induced ZZ coupling (there are no victim gates
to induce it) and must therefore be attributed to the readout chain.

In practice, the control is realized as follows:

\begin{enumerate}[leftmargin=*]
    \item \textbf{Qubit allocation:} Victim qubits $\mathcal{V}$ and attacker qubits
    $\mathcal{A}$ are allocated on the QPU identically to the gate-crosstalk experiments.
    \item \textbf{Victim circuit:} Qubits in $\mathcal{V}$ are initialised to the
    computational ground state $|0\rangle^{\otimes n_V}$ and never operated on---zero
    gates, zero microwave pulses.
    \item \textbf{Attacker circuit:} The attacker probe circuit executes identically to
    the gate-crosstalk experiments (idle or single-CNOT probe, same shot count).
    \item \textbf{Measurement:} Only the attacker qubits $\mathcal{A}$ are measured,
    identical to the gate-crosstalk experiments.
    \item \textbf{Leakage quantification:} TVD between the attacker's measurement
    distribution in this configuration and the attacker-alone (victim absent)
    configuration quantifies the readout-only leakage baseline.
\end{enumerate}

The disambiguation works because the two crosstalk mechanisms produce
\textbf{qualitatively different distance-decay signatures} (Section~5.3 of the revised
manuscript): gate-induced ZZ coupling decays exponentially with SWAP-path distance
($d_0 \approx 1$--$2$ hops), while readout crosstalk---mediated by the
frequency-multiplexed measurement feedline, JPA chain, and ADC paths---is approximately
distance-independent ($d_0 \gg 10$ hops).

\medskip\noindent\textbf{Changes to manuscript:} We have added a more detailed
step-by-step description of the control realisation in Section~5.3 and clarified that the
control is an experimental methodology, not part of the attack protocol. We have also
added explicit language distinguishing ``measurement-induced readout crosstalk'' (the
physical effect) from the ``readout-only control'' (the experimental technique used to
isolate it).

% ── R1C2 ─────────────────────────────────────────────────────────
\subsection*{Comment 2}

\noindent\textbf{\textit{``Now that you use the measurement-induced crosstalk, does it
mean that your protocol is in an indeterministic way?''}}

\medskip\noindent\textbf{Response:}

No, the protocol is not indeterministic. We believe there may be a terminological
misunderstanding that we have clarified in the revised manuscript.

The readout-only control is used to \textbf{separate and quantify} measurement-induced
crosstalk as a background effect, not to rely on it. The forward predictive model itself
(the CrossEcho protocol) predicts gate-induced ZZ crosstalk leakage from device calibration
data---this prediction is deterministic given the inputs (qubit frequencies, coupling
strengths, circuit parameters). The Monte Carlo component propagates \textbf{epistemic
uncertainty} (uncertainty in the input parameters due to limited calibration data), not
intrinsic quantum randomness.

The attacker's distinguishing game (Section~3) does involve quantum measurement randomness,
but this is fully characterised by the shot-count statistics and the TVD framework. The
accumulated ZZ phase itself is a deterministic consequence of the Hamiltonian evolution;
the randomness enters only through the projective measurement at the end. This is no more
``indeterministic'' than any other quantum measurement---and the same statistical framework
(TVD, Neyman--Pearson, Chernoff bounds) applies.

To be explicit about the distinction:
\begin{itemize}[leftmargin=*]
    \item \textbf{Gate-induced ZZ crosstalk} (deterministic phase accumulation
    $\rightarrow$ measured via population statistics)
    \item \textbf{Readout crosstalk} (classical/microwave parasitic coupling $\rightarrow$
    measured as a distance-independent baseline)
    \item \textbf{Monte Carlo propagation} (epistemic uncertainty quantification, not
    quantum indeterminism)
\end{itemize}

\medskip\noindent\textbf{Changes to manuscript:} We have added a clarification paragraph
in Section~5.3 explicitly distinguishing deterministic phase accumulation from measurement
randomness, and added language in Section~4.4 clarifying that the Monte Carlo propagation
addresses epistemic (parameter) uncertainty, not intrinsic quantum indeterminism.

% ── R1C3 ─────────────────────────────────────────────────────────
\subsection*{Comment 3}

\noindent\textbf{\textit{``You claimed that `Cross-architecture validation on the Origin
Quantum Cloud WK\_C180\_2 processor reveals an order-of-magnitude increase in crosstalk
magnitude relative to IBM, directly confirming the architecture-dependence that the model
identifies as its dominant uncertainty source'. What is the cross-architecture? Is it
related to your protocol or quantum processors? If the previous one, it is not appropriate
to use two different quantum hardware, i.e.\ Origin Quantum Cloud WK\_C180\_2 processor
and IBM Quantum Heron r2 processors, to validate this conclusion.''}}

\medskip\noindent\textbf{Response:}

We appreciate this comment, which identifies an ambiguity in our terminology.
``Cross-architecture'' in our manuscript refers to \textbf{different quantum processor
architectures}---specifically, IBM's heavy-hexagonal topology (used in Heron~r2) versus
Origin's 2D grid topology (used in WK\_C180\_2). It is not a claim about the CrossEcho
protocol itself being ``cross-architecture'' in some protocol-level sense.

The logic of our validation is as follows:

\begin{enumerate}[leftmargin=*]
    \item The perturbation-theory model (Eqs.~4--6) predicts that crosstalk magnitude
    depends on qubit frequencies ($\omega_i$), coupling strengths ($g_{ij}$), and lattice
    topology---all of which vary across processor architectures.
    \item If crosstalk were dominated by readout crosstalk alone (which is
    architecture-independent, being mediated by the classical measurement chain), we would
    expect similar TVD magnitudes on any superconducting processor with comparable readout
    fidelity, regardless of qubit topology.
    \item The observation of an order-of-magnitude difference between IBM Heron~r2 and
    Origin WK\_C180\_2 directly \textbf{falsifies} the architecture-independence
    hypothesis and supports the architecture-dependence prediction of the
    perturbation-theory model.
\end{enumerate}

Using two genuinely different hardware platforms is precisely the \textbf{correct} way to
test an architecture-dependence claim---testing on only one architecture would leave the
claim untested.

\medskip\noindent\textbf{Changes to manuscript:} We have rewritten the relevant passages
to replace ``cross-architecture'' with explicit language: ``cross-processor-architecture
validation'' or ``architecture comparison between IBM (heavy-hex) and Origin (2D grid)
topologies.'' We have also added a sentence in Section~6.4 explaining why testing on two
different architectures is the appropriate validation strategy for the
architecture-dependence hypothesis.

% ── R1C4 ─────────────────────────────────────────────────────────
\subsection*{Comment 4}

\noindent\textbf{\textit{``What is the Origin Quantum Cloud WK\_C180\_2 processor?''}}

\medskip\noindent\textbf{Response:}

The Origin Quantum Cloud WK\_C180\_2 is a 180-qubit superconducting transmon processor
accessible through the Origin Quantum Cloud platform (\url{https://qcloud.originqc.com.cn}).
We have expanded its description in the revised manuscript. Key specifications:

\begin{itemize}[leftmargin=*]
    \item \textbf{Qubit count:} 180 superconducting transmon qubits
    \item \textbf{Topology:} 2D rectangular grid (as opposed to IBM's heavy-hexagonal
    lattice)
    \item \textbf{Qubit type:} Fixed-frequency transmons (no tunable couplers)
    \item \textbf{Native gate set:} CZ (controlled-Z) and RPhi (single-qubit rotation
    with arbitrary phase)
    \item \textbf{Readout:} Dispersive readout via frequency-multiplexed resonators on a
    common feedline
    \item \textbf{Access:} Python SDK \texttt{pyqpanda3} v0.3.5
    \item \textbf{Calibration data exposed:} Limited compared to IBM; qubit frequencies
    and gate fidelities are partially available through the API
\end{itemize}

We selected WK\_C180\_2 for architectural contrast with IBM's Heron~r2 because:
(i)~the denser 2D-grid connectivity produces larger baseline capacitive coupling between
qubits, (ii)~the fixed-frequency design (without tunable couplers) provides less dynamical
crosstalk suppression than IBM's tunable-coupler architecture, and (iii)~the estimated
inter-qubit detuning ($\Delta \approx 200$--$300$~MHz) is smaller than IBM's engineered
$\Delta \gtrsim 500$~MHz, producing larger ZZ coupling via the $\Delta^{-2}$ dependence in
Eq.~4.

\medskip\noindent\textbf{Changes to manuscript:} We have added a dedicated paragraph
describing the WK\_C180\_2 processor and the rationale for its selection as a
cross-processor-architecture comparison platform in Section~5.2, along with a full
comparison table (Table~2) contrasting IBM Heron~r2 and Origin WK\_C180\_2 across
topology, coupler type, gate set, estimated ZZ coupling, SDK, and calibration data
availability.

% ── R1C5 ─────────────────────────────────────────────────────────
\subsection*{Comment 5}

\noindent\textbf{\textit{``You introduced a terminology of `crossecho', and does it have
something to do with the similar concept `echo' in quantum communication?''}}

\medskip\noindent\textbf{Response:}

We thank the reviewer for this observation. The name ``CrossEcho'' is \textbf{not}
directly related to the ``echo'' concept in quantum communication (Hahn spin echoes,
dynamical decoupling, or echo-based entanglement purification protocols). We recognise
that the term ``echo'' carries established meaning in the quantum information community,
and we have clarified the intended meaning in the revised manuscript.

The ``echo'' in CrossEcho is a descriptive metaphor for the physical mechanism:
\begin{itemize}[leftmargin=*]
    \item \textbf{``Cross''} refers to \textbf{crosstalk}---the unwanted inter-qubit
    coupling that propagates through the coupling fabric.
    \item \textbf{``Echo''} refers to the fact that a victim's gate operations leave an
    \textbf{imprint} on neighbouring attacker qubits through ZZ phase accumulation. Just
    as an acoustic echo is a delayed copy of a sound reflecting off a surface, the ZZ
    phase on the attacker's qubits is a ``copy'' of the victim's gate activity propagating
    through the coupling graph.
\end{itemize}

Physically, the mechanism is \textbf{phase accumulation}
($\Phi_{\mathrm{acc}} = 2\pi|\zeta_{\mathrm{eff}}|\,t_{\mathrm{gate}}\,\rho_{\mathrm{2Q}}\,D\,n_V\,s$,
Eq.~9), not population transfer (which would be a SWAP or iSWAP interaction). The
accumulated phase is detected as a change in the attacker's measurement statistics
(population difference after state tomography).

There is a loose analogy with spin echoes: both involve controlled phase evolution. In a
Hahn echo, a $\pi$-pulse refocuses dephasing by reversing the sign of the
system--environment interaction. In CrossEcho, the ZZ coupling imprints a phase on the
attacker qubits that is proportional to the victim's two-qubit gate count. However, the
mechanisms are fundamentally different---CrossEcho does not involve any echo pulse sequence
or refocusing.

\medskip\noindent\textbf{Changes to manuscript:} We have added a clear definition of the
``echo'' terminology in Section~1 (Introduction), explicitly stating that it refers to ZZ
phase accumulation (not population transfer) imprinted through the coupling fabric, and
distinguishing it from the Hahn/spin-echo concept in quantum communication. This definition
is cross-referenced in the abstract.

% ===========================================================================
\section*{Reviewer 2}
% ===========================================================================

% ── R2C1 ─────────────────────────────────────────────────────────
\subsection*{Comment 1}

\noindent\textbf{\textit{``The Monte Carlo sampling assumes normal distributions for
$g_{ij}$ and $\Delta_{ij}$. Justify this Gaussian assumption, or briefly discuss how
non-Gaussian fabrication variability would alter uncertainty bounds.''}}

\medskip\noindent\textbf{Response:}

We thank the reviewer for this methodological question. The Gaussian assumption is a
pragmatic first-order choice, and we have added a justification and a discussion of
non-Gaussian effects in the revised manuscript (Section~4.4).

\medskip\noindent\textbf{Justification of the Gaussian assumption:}

\begin{enumerate}[leftmargin=*]
    \item \textbf{Fabrication physics:} Josephson junction resistance ($R_n$), which sets
    the Josephson energy $E_J$ and therefore the qubit frequency $\omega_q \propto
    \sqrt{E_J E_C}$, is determined by the oxidation tunnel barrier thickness. For nominally
    identical junctions on the same wafer, barrier thickness variations have been
    experimentally characterised to follow approximately normal distributions
    (Osman~\emph{et al.}, 2023; Alghadeer~\emph{et al.}, 2025).
    \item \textbf{Central Limit Theorem:} The qubit frequency $\omega_q$ and coupling
    strength $g_{ij}$ depend on multiple independent fabrication steps (junction oxidation,
    capacitor lithography, resonator definition). By the Central Limit Theorem, the
    aggregate distribution tends toward Gaussian even if individual step distributions
    are not.
    \item \textbf{Maximum entropy:} In the absence of detailed foundry distribution data
    (which is proprietary to IBM and Origin), the Gaussian distribution is the
    maximum-entropy (least-informative) distribution for a given mean and variance. Any
    other choice would implicitly encode additional assumptions not supported by available
    data.
\end{enumerate}

\medskip\noindent\textbf{Effect of non-Gaussian variability:}

If the true fabrication distribution deviates from Gaussian, our uncertainty bounds would
be affected as follows:

\begin{itemize}[leftmargin=*]
    \item \textbf{Heavy-tailed distributions} (e.g., log-normal junction resistances):
    The tails of the predicted TVD distribution would be wider than our $2\sigma$ Gaussian
    intervals, meaning our confidence intervals would be \textbf{anti-conservative} (they
    would understate the true uncertainty). The mean prediction $L_{\mathrm{pred}}$ could
    also be biased if the distribution is skewed.
    \item \textbf{Bimodal distributions} (e.g., two fabrication batches with different
    mean parameters): The mean prediction could be misleading, falling between two
    clusters. In this case, the model should be stratified by fabrication batch.
    \item \textbf{Bounded uniform distributions} (e.g., if $\Delta$ is constrained by
    design to $[500, 700]$~MHz): Our Gaussian intervals would be \textbf{conservative}
    (they would overstate the uncertainty), since the true distribution would have lighter
    tails than Gaussian.
\end{itemize}

For the present work, the dominant error is not the shape of the input distribution but the
systematic offset in its mean: $\bar{g}$ and $\bar{\Delta}$ are estimated without access to
actual qubit frequencies. As noted in Section~7.1, frequency opaqueness contributes
$\sim$90\% of the model error budget, dwarfing any distribution-shape effects. Future work
with access to coarse frequency bins from providers would enable empirical calibration of
the input distributions.

\medskip\noindent\textbf{Changes to manuscript:} We have added a paragraph in Section~4.4
justifying the Gaussian assumption (three grounds: fabrication physics, CLT, maximum
entropy) and a paragraph discussing the effect of non-Gaussian variability on the
uncertainty bounds (heavy-tailed, bimodal, and bounded-uniform cases).

% ── R2C2 ─────────────────────────────────────────────────────────
\subsection*{Comment 2}

\noindent\textbf{\textit{``Figure~4 shows non-monotonic TVD fluctuations across SWAP
distances for both gate and readout crosstalk. Add a short discussion on the origins of
this noise (binomial sampling error, per-qubit readout fidelity variation) and how
increased shot counts smooth this distance profile.''}}

\medskip\noindent\textbf{Response:}

We agree that the non-monotonic fluctuations in Figure~4 (now Figure~2 in the revised
manuscript) merit explicit discussion. We have added a paragraph in Section~6.2 addressing
the three noise sources and their shot-count scaling.

\medskip\noindent\textbf{Origins of the non-monotonic fluctuations:}

\begin{enumerate}[leftmargin=*]
    \item \textbf{Binomial sampling error:} At $N = 2000$ shots, the standard error of a
    single TVD measurement is $\sigma_{\mathrm{TVD}} \approx 1/\sqrt{N} \approx 0.022$ in
    the worst case (uniform distribution). This is comparable to the measured TVD values
    themselves ($0.005$--$0.057$), meaning individual cell measurements are noise-dominated.
    This is the primary source of the fluctuations.
    \item \textbf{Per-qubit readout fidelity variation:} Different qubits on the same
    processor have different readout error rates (typically $1$--$5\%$ on IBM Heron~r2).
    A cell placed on qubits with higher readout error will show higher apparent TVD even if
    the true crosstalk is identical, producing cell-to-cell scatter at the same distance.
    \item \textbf{Per-path coherence variation:} Two paths with the same hop-count $d$ may
    traverse intermediate qubits with different $T_1$ and $T_2$ values, producing different
    effective $\zeta_{\mathrm{eff}}$ via decoherence during SWAP-path propagation.
\end{enumerate}

\medskip\noindent\textbf{Effect of increased shot counts:}

The binomial sampling error scales as $\propto 1/\sqrt{N}$. Increasing from $N = 2000$ to
$N = 20{,}000$ would reduce the per-cell standard error from $\sim 0.022$ to $\sim 0.007$,
and to $N = 10^5$ would reduce it to $\sim 0.003$. At $10^5$ shots, the binomial noise
would be well below the typical measured TVD values, and the underlying distance-decay
profile would emerge clearly.

Per-qubit readout fidelity variation could be partially corrected by normalising each
cell's TVD by the geometric mean of its qubits' readout error rates:
$\mathrm{TVD}_{\mathrm{norm}} = \mathrm{TVD}_{\mathrm{meas}} /
\sqrt{\epsilon_{\mathrm{ro},q_1} \epsilon_{\mathrm{ro},q_2}}$.
This is straightforward to implement (readout error rates are exposed by the IBM API) and
would remove one source of spurious cell-to-cell variation. We did not apply this
correction in the present analysis to keep the pipeline simple, but we note it as an
available refinement for future studies.

\medskip\noindent\textbf{Changes to manuscript:} We have added a dedicated subsection in
Section~6.2 discussing the three noise sources, their relative magnitudes, and their
shot-count scaling. We have also added a sentence noting the readout-normalisation option.

% ── R2C3 ─────────────────────────────────────────────────────────
\subsection*{Comment 3}

\noindent\textbf{\textit{``Table~5 compares TVD metrics against classification accuracy
from prior attack papers, but does not convert earlier accuracy results into equivalent TVD
bounds for direct side-by-side comparison. Add a brief supplementary calculation to unify
cross-study leakage quantification.''}}

\medskip\noindent\textbf{Response:}

This is an excellent suggestion. We have added a supplementary calculation that converts
prior accuracy-based metrics into equivalent TVD lower bounds, enabling direct quantitative
comparison. The calculation is included in Section~6.5 and summarised in the revised
Table~5.

\medskip\noindent\textbf{Method:} For a binary hypothesis test with $N$ independent
measurements, classification accuracy $A$ is bounded above by the total variation distance
via the Neyman--Pearson lemma and Chernoff bound:
%
\begin{equation}
    A \leq \frac{1}{2} + \frac{1}{2} \min(1, N \cdot \mathrm{TVD}).
    \label{eq:accuracy_bound}
\end{equation}
%
Rearranging, the minimum TVD required to achieve accuracy $A$ (a lower bound on the actual
TVD) is:
%
\begin{equation}
    \mathrm{TVD}_{\mathrm{min}} \geq \max\left(0, \frac{2A - 1}{N}\right).
    \label{eq:tvd_lower_bound}
\end{equation}

\medskip\noindent\textbf{Application to prior studies:}
Hu~\emph{et al.} (2025) report $100\%$ circuit-size prediction accuracy with typical
$N \approx 2000$--$4000$ shots, yielding $\mathrm{TVD}_{\mathrm{min}} \geq 0.00025$--$0.0005$.
Choudhury~\emph{et al.} (2024) report $85.7\%$ algorithm classification accuracy with
$N \approx 4000$, yielding $\mathrm{TVD}_{\mathrm{min}} \geq 0.00018$.

Our measured TVD range of $0.005$--$0.057$ is \textbf{10--300$\times$ above these lower
bounds}, confirming consistency: the TVD values we measure are more than sufficient to
support the classification accuracies reported in prior work. We note that these are
\textbf{lower bounds}; actual TVD values in prior studies are likely higher because
(i)~classifiers optimise over multiple measurements, (ii)~the Chernoff bound is loose for
small $N$, and (iii)~multi-class classification accuracy yields a looser binary bound than
a dedicated binary hypothesis test.

\medskip\noindent\textbf{Changes to manuscript:} We have added the TVD--accuracy conversion
method (Eqs.~15--16) as a paragraph in Section~6.5 and expanded Table~5 to include a
rightmost column with the equivalent TVD lower bounds for each prior study.

% ── R2C4 ─────────────────────────────────────────────────────────
\subsection*{Comment 4}

\noindent\textbf{\textit{``Standardize shorthands on first use: ZZ coupling, TVD, NISQ,
QPU, RMSE, SWAP-path distance, JPA, ADC.''}}

\medskip\noindent\textbf{Response:}

We thank the reviewer for this careful check. We have performed a full pass over the
manuscript to ensure all abbreviations are defined on first use in the body text. The
\verb|\abbr{}| block on the title page has been expanded, and we have verified each
abbreviation's first-use location.

Specifically:
\begin{itemize}[leftmargin=*]
    \item \textbf{ZZ coupling:} Defined in the abstract and again in Section~2.1 (first
    technical use as ``static ZZ coupling'' with the physical mechanism described)
    \item \textbf{TVD (Total Variation Distance):} Defined in the abstract and
    Section~3.2 (Eq.~2)
    \item \textbf{NISQ (Noisy Intermediate-Scale Quantum):} Defined in the abstract and
    spelled out at first body-text use in Section~1: ``Preskill's Noisy Intermediate-Scale
    Quantum (NISQ) paradigm''
    \item \textbf{QPU (Quantum Processing Unit):} Spelled out at first body-text use in
    Section~1: ``quantum processing units (QPUs)''
    \item \textbf{RMSE (Root-Mean-Square Error):} Spelled out at first technical use in
    Section~5.4: ``root-mean-square error (RMSE)''
    \item \textbf{SWAP-path distance:} Defined in Section~4.2: ``SWAP-mediated exchange
    interactions---what we refer to as the SWAP path (exchange interaction path)''
    \item \textbf{JPA (Josephson Parametric Amplifier):} Spelled out at first use in
    Section~2.3
    \item \textbf{ADC (Analogue-to-Digital Converter):} Spelled out at first use in
    Section~2.3
\end{itemize}

We have also verified that all abbreviations used in figures and tables are defined in the
surrounding text or in the expanded \verb|\abbr{}| block (now including ECR, CP, CZ, CNOT).

\medskip\noindent\textbf{Changes to manuscript:} We have added missing first-use
definitions (NISQ, CNOT, ECR, RMSE, SWAP-path) in the body text, expanded the
\verb|\abbr{}| block to be comprehensive (adding JPA, ADC, CP, CZ, CNOT), and verified
every abbreviation used in the manuscript has a single, clear first-use definition.

% ── R2C5 ─────────────────────────────────────────────────────────
\subsection*{Comment 5}

\noindent\textbf{\textit{``Remove redundant repetitive descriptions of the readout-only
control experiment across Section~5, 6 and 7; consolidate core methodology once and
cross-reference later sections.''}}

\medskip\noindent\textbf{Response:}

We agree that the original manuscript contained redundant descriptions. The readout-only
control was described in three places: Section~5.3 (methods), Section~6.2 (results), and
Section~7 (discussion). We have now consolidated as follows:

\begin{itemize}[leftmargin=*]
    \item \textbf{Section~5.3 (Methods):} Contains the \textbf{complete, authoritative
    description} of the readout-only control---rationale, detailed five-step procedure,
    physical mechanisms isolated, and expected distance-decay signatures. This is the
    single source of truth.
    \item \textbf{Section~6.2 (Results):} Reports the findings from the control experiment
    (flat distance decay, $d_0 \gg 10$, statistical tests) with a \textbf{single-sentence
    cross-reference} to Section~5.3 for the methodology, rather than re-describing it.
    \item \textbf{Section~7 (Discussion):} Discusses implications of the control results
    for prior work and future methodology, \textbf{cross-referencing} Section~5.3 rather
    than re-describing the experiment.
\end{itemize}

This follows the standard IMRaD structure: Methods describe what was done, Results report
what was found, Discussion interprets what it means.

\medskip\noindent\textbf{Changes to manuscript:} We have removed the redundant methodology
descriptions from Sections~6.2 and 7, replacing each with a one-sentence cross-reference
to Section~5.3, and verified that the consolidated description in Section~5.3 is complete
and self-contained.

% ── R2C6 ─────────────────────────────────────────────────────────
\subsection*{Comment 6}

\noindent\textbf{\textit{``The term `echo' in `CrossEcho' is evocative but not rigorously
defined---clarify whether it refers to phase accumulation, population transfer, or
both.''}}

\medskip\noindent\textbf{Response:}

This comment aligns with Reviewer~1's Comment~5, and we have addressed both through a
unified clarification in the revised manuscript.

The ``echo'' in CrossEcho refers specifically to \textbf{phase accumulation}---the ZZ phase
$\Phi_{\mathrm{acc}}$ imprinted on the attacker's qubits by the victim's two-qubit gate
operations, propagating through the coupling fabric. It is \textbf{not} population transfer
(that would require a resonant SWAP or iSWAP interaction, which is a different coupling
mechanism). The accumulated phase is detected as a change in the attacker's measurement
statistics (population difference) after computational-basis measurement.

We have added the following definition in Section~1:
\begin{quote}
\small
``The project name CrossEcho reflects this mechanism: `Cross' refers to the crosstalk
propagating across the coupling topology, and `Echo' refers specifically to the ZZ phase
accumulated on the attacker's qubits ($\Phi_{\mathrm{acc}}$ in Eq.~9) in proportion to the
victim's two-qubit gate count. This is phase accumulation, not population transfer (which
would require resonant SWAP or iSWAP interactions), and is distinct from the Hahn
spin-echo or dynamical-decoupling `echo' concepts in quantum communication and NMR.''
\end{quote}

\medskip\noindent\textbf{Changes to manuscript:} We have added a precise definition of
``echo'' in Section~1, explicitly stating it refers to ZZ phase accumulation (not
population transfer) and distinguishing it from Hahn echoes. This is cross-referenced in
the abstract.

% ── R2C7 ─────────────────────────────────────────────────────────
\subsection*{Comment 7}

\noindent\textbf{\textit{``Compared with IBM Quantum Heron r2 processors, what are the
differences of Origin Quantum Cloud? What are the respective advantages and disadvantages
of each quantum cloud?''}}

\medskip\noindent\textbf{Response:}

We thank the reviewer for this question, which prompted us to add a substantive comparison
that strengthens the cross-architecture analysis. We have added a comparison table and
discussion in Section~5.2.

\medskip\noindent\textbf{Summary of differences:}

\begin{table}[h]
\centering
\caption{Comparison of IBM Heron~r2 and Origin WK\_C180\_2 quantum processor architectures.}
\label{tab:platform}
\small
\begin{tabular}{p{3.2cm}p{2.5cm}p{2.5cm}}
\toprule
Feature & IBM Heron~r2 & Origin WK\_C180\_2 \\
\midrule
Qubit count & 156 & 180 \\
Topology & Heavy-hexagonal & 2D rectangular grid \\
Connectivity & 2--3 neighbours/qubit & 4 neighbours/qubit (bulk) \\
Coupler type & Tunable & Fixed \\
Native 2Q gate & ECR & CZ $+$ RPhi \\
Frequency engineering & $\Delta \gtrsim 500$~MHz & $\Delta \approx 200$--$300$~MHz \\
Typical ZZ coupling & $\sim 10$--$50$~kHz & $\sim 100$--$500$~kHz \\
SDK & Qiskit Runtime & pyqpanda3 \\
Calibration data & T1, T2, gate/readout errors, coupling map & Limited (partial) \\
User community & Large, well-documented & Smaller, growing \\
\bottomrule
\end{tabular}
\end{table}

\medskip\noindent\textbf{Advantages of IBM Heron~r2:}
\begin{itemize}[leftmargin=*]
    \item Tunable couplers enable dynamical ZZ suppression, reducing baseline crosstalk
    \item Engineered frequency detuning ($\Delta \gtrsim 500$~MHz) further suppresses ZZ
    via the $\Delta^{-2}$ dependence
    \item Mature SDK (Qiskit) with extensive documentation and community support
    \item Rich calibration data exposed through the API enables detailed noise modelling
    \item Heavy-hex topology supports efficient surface-code error correction
\end{itemize}

\noindent\textbf{Disadvantages of IBM Heron~r2:}
\begin{itemize}[leftmargin=*]
    \item Heavy-hex connectivity limits qubit allocation flexibility
    \item QPU time is competitive and queue times can be long during peak usage
    \item Frequency opaqueness limits independent crosstalk modelling (Section~7.1)
\end{itemize}

\noindent\textbf{Advantages of Origin WK\_C180\_2:}
\begin{itemize}[leftmargin=*]
    \item Denser 2D-grid connectivity provides more qubit allocation options
    \item Academic access programme provides dedicated QPU time
    \item Native CZ gate offers an alternative entangling primitive
\end{itemize}

\noindent\textbf{Disadvantages of Origin WK\_C180\_2:}
\begin{itemize}[leftmargin=*]
    \item Fixed-frequency transmons without tunable couplers $\rightarrow$ larger baseline
    $g$, higher ZZ crosstalk
    \item Smaller inter-qubit detuning $\rightarrow$ larger $|\zeta|$ via the $\Delta^{-2}$
    dependence
    \item Limited calibration data exposed through the API restricts independent noise
    characterisation
    \item Smaller user community and less mature SDK
\end{itemize}

\medskip\noindent\textbf{Relevance to our study:} The $15\times$ CNOT TVD excess on Origin
relative to IBM is a direct consequence of these architectural differences: denser
connectivity (larger $g$) and smaller detuning (smaller $\Delta$) combine via Eq.~4 to
produce $\sim$10--$25\times$ larger $|\zeta|$. This validates the perturbation-theory
model's central claim: crosstalk magnitude is strongly architecture-dependent.

\medskip\noindent\textbf{Changes to manuscript:} We have added a comparison table and
discussion in Section~5.2, and cross-referenced this comparison in the cross-processor
results discussion (Section~6.4).

% ===========================================================================
\section*{Additional Changes}
% ===========================================================================

In addition to the reviewer-requested changes, we have made the following improvements:

\begin{enumerate}[leftmargin=*]
    \item \textbf{Updated the \texttt{\textbackslash abbr\{\}} block} to be comprehensive,
    including all abbreviations used in the manuscript (now: NISQ, QPU, TVD, ZZ, QPT,
    SWAP, RMSE, ECR, JPA, ADC, CP, CZ, CNOT).
    \item \textbf{Verified all references} against Crossref and arXiv, correcting any
    inconsistencies.
    \item \textbf{Updated the revision date} in the manuscript metadata.
    \item \textbf{Standardised figure cross-references} to use the correct numbering after
    consolidation.
\end{enumerate}

\vspace{12pt}
\noindent
We believe these revisions have substantially improved the clarity, rigour, and
completeness of the manuscript. We thank both reviewers again for their time and
constructive input.

\vspace{24pt}
\noindent\textbf{Corresponding Author:}\\
Justice Owusu Agyemang\\
Quantum and Assistive Technologies Lab\\
Kwame Nkrumah University of Science and Technology\\
\texttt{justice.agyemang@knust.edu.gh}

\end{document}
